\section{Introduction}
\label{sec:1}

Extended reality offers a shared environment in which people can communicate through embodied representations and interact with common task materials; research has examined such environments for group discussion, collaborative learning, and medical team simulation ~\cite{p70,p51,p66}. These environments also support observation of gaze, movement, object interactions, and task events, creating opportunities to connect multimodal sensing with assistance ~\cite{p34,p57,pro2}. The starting point of this work is therefore an established sensing and interaction capability, rather than an assumption that XR automatically improves teamwork.

A first line of research connects multimodal sensing to individual assistance. Pro\textsuperscript{2}Assist follows procedural activity through AR glasses and provides step-aware guidance ~\cite{pro2}. GuideAI adapts learning support using physiological and behavioral information, and AwareLLM combines physiological and contextual sensing for personalized knowledge-work assistance ~\cite{guide,aware}. Sensing What Surveys Miss uses electrodermal activity and mouse behavior to time proactive LLM support, demonstrating that assistance timing can be evaluated against alternative timing conditions ~\cite{sensing}. These studies establish substantial precedent for state-aware assistance; physiological sensing coupled with an LLM is not, by itself, the contribution proposed here.

A second line moves from sensing an individual to examining relationships among collaborators. Physiological and eye-movement coupling has been used to assess remote collaborative processes, while shared arousal events have been aligned with collaborative problem solving and regulation ~\cite{p8,p18,p47}. In a three-person VR task, teammates' physiological and behavioral information was used to predict another member's upcoming actions, relating predictability to team performance ~\cite{p52}. Recent medical-team research further combines physiological synchrony with dialogue analysis, showing that synchrony peaks can accompany both shared discovery and shared uncertainty ~\cite{huang}. Thus, collaborative sensing provides potentially useful evidence, but interpreting its social meaning requires behavioral and task context ~\cite{p32,p56}.

The missing connection addressed here is the decision from relational evidence to assistance. Studies of social biofeedback primarily investigate how people communicate or interpret shared bodily information, whereas studies of team physiology investigate associations, temporal patterns, or prediction ~\cite{p38,p46,p49,p52}. XR moderation and attention guidance provide intervention mechanisms, but their findings do not establish the incremental benefit of adding multi-person physiology to an otherwise comparable agent ~\cite{p34,p70}. Taken together, these research strands leave an empirical question: under what conditions does multi-person physiology change an agent's target, timing, or action in ways that improve collaboration beyond behavioral and task sensing alone?

Consider a team inspecting different pieces of evidence in a shared XR workspace. A member's silence could warrant an invitation if important information remains unexpressed, but an invitation could be unnecessary if that member is still inspecting the material. Prior work on multimodal collaboration measurement and proactive interaction motivates evaluating such interpretations against observable task events, participant accounts, and the consequences of an intervention ~\cite{p57,speaker,expect}. The proposed agent therefore estimates an actionable support need rather than presenting an inferred emotion as fact.

We propose to study a constrained set of decisions: continue observing, privately ask whether help is wanted, invite a contribution, prompt a shared check, or defer an episode to reflection. These are proposed design choices motivated by existing moderation, pedagogical scaffolding, and human-AI interaction work ~\cite{p70,p27,p3}. The primary study will examine one collaborative XR task and one operationally defined support need. Medical training and education will extend this mechanism only after its measurement and intervention logic have been validated. This sequencing allows the core mechanism to be validated before domain-specific adaptation; differences in medical teamwork assessment and collaborative learning regulation motivate separate outcome definitions ~\cite{p66,p29,p51}.

The intended contributions are: (1) a task-grounded representation of team support needs that distinguishes evidence, uncertainty, and action; (2) a relational sensing and intervention pipeline whose individual components can be compared with simpler alternatives; and (3) controlled evidence about when multi-person physiology adds value, fails to add value, or increases intervention cost. For training and education, the evaluation additionally tests behavior and understanding after assistance is removed, because performance during support is insufficient evidence of durable learning ~\cite{p62}.
